\documentclass[10pt,a4paper]{amsart}
\usepackage[margin=19mm]{geometry}
\usepackage{amsmath,amssymb,mathtools}
\usepackage[scr=rsfs,cal=euler]{mathalfa}
\usepackage{xfrac}
\usepackage[unicode,hidelinks,bookmarks=false]{hyperref}
\newcommand{\ie}{i.e.\ }
\newcommand{\cf}{cf.\ }
\newcommand{\resp}{respectively\ }
\newcommand{\Subsection}{\subsection}
\providecommand{\nobreakdash}{\nobreak\hbox{-}}
\setlength{\emergencystretch}{2em}
\multlinegap0pt
\usepackage{bm}
\usepackage{bbm}
\usepackage{dsfont}
\usepackage{xspace}
\usepackage{cancel}
\usepackage{mathtools}
\usepackage{amsthm}
\makeatletter
\@ifundefined{theorem}{\newtheorem{theorem}{Theorem}}{}
\@ifundefined{lemma}{\newtheorem{lemma}[theorem]{Lemma}}{}
\@ifundefined{proposition}{\newtheorem{proposition}[theorem]{Proposition}}{}
\makeatother
\makeatletter
\DeclareMathOperator{\diam}{diam}
\expandafter\ifx\csname proof\endcsname\relax
\renewenvironment{proof}[1][\proofname]
{\par\pushQED{\qed}
	\normalfont\topsep6\p@\@plus6\p@\relax\trivlist
	\item[\hskip\labelsep\bfseries#1\@addpunct{.}]
	\ignorespaces}
\fi
\makeatother
\title[Asymptotically attaining the Moore bound]{Asymptotically attaining the Moore bound}
\author{Wouter Cames van Batenburg, Samuel Korsky}
\date{Source: arXiv:2608.03965v1 (CC BY 4.0)}
\begin{document}
\maketitle

\noindent\emph{Source and licence.} Asymptotically attaining the Moore bound. Version arXiv:2608.03965v1 (CC BY 4.0), distributed under the Creative Commons Attribution 4.0 International licence, as recorded in the arXiv licence metadata for this exact version and shown on the abstract page https://arxiv.org/abs/2608.03965v1. Full rights notice: licenses/arxiv-2608.03965-CC-BY-4.0.txt.\par\smallskip

\noindent\textit{Editorial scope.} Counted unit: the Proposition whose source label is \texttt{prop:H} in the article (source0.tex lines 254--272, source label \texttt{prop:H}), delivered together with the article's own complete original proof of it (source0.tex lines 274--334, one \texttt{proof} environment of 61 lines). No mathematics is written, repaired or re-derived: statement and proof are the article's own text line for line. Required context retained uncounted: lem:routing (lines 188--192). Every definition, display and label of those lines is retained unchanged. Omitted from this package and not redistributed: 1--187, 193--253, 273--273, 335--563. Nothing inside a retained excerpt is omitted or altered: every retained body is delivered exactly as the article's lines are written, so no omission receipt of any kind accompanies this package. Citations: the retained excerpts carry no citation command at all, so no bibliography entry of the article is transcribed and no reference is redistributed. Reference adaptation: none was needed and none was made; every reference inside the delivered unit resolves inside it, so no unresolved reference or citation is produced. Printed slips kept literal: no printed word, symbol, hypothesis or number of the counted statement or of its proof was altered. Layout and class: the article's own class is \texttt{article} with packages including inputenc, fontenc, mathpazo, comment, babel, geometry, parskip, xcolor, enumitem, amsmath, amsfonts, amssymb, amsthm, mathtools; the wrapper is the lane's pinned \texttt{amsart} editorial wrapper at 10pt on A4, which restates the theorem-like environments the retained text needs and adds the article's own macro definitions that the retained text needs (\texttt{\textbackslash diam}), with extra wrapper packages (bm, bbm, dsfont, xspace, cancel, mathtools). The delivered numbering is the unit's own. Assets: no retained span contains a figure environment, a graphics-inclusion command, a PDF-page inclusion, a table or a TikZ picture; no image file is copied into this package, the package declares no asset dependency and no image file is redistributed with it. Licence: version arXiv:2608.03965v1 of this article is distributed under the Creative Commons Attribution 4.0 International licence, as recorded in the arXiv licence metadata for this exact version and shown on the abstract page https://arxiv.org/abs/2608.03965v1. A full rights notice is delivered at licenses/arxiv-2608.03965-CC-BY-4.0.txt. Characters: every retained excerpt is pure ASCII, so the delivered engine, XeLaTeX with its default Latin Modern text font, reports no missing character.
\begin{lemma}\label{lem:routing}
Let $F$ and $F'$ be complete flags in a $t$-dimensional vector space.  There is a sequence of complete flags
$$F=G^0,G^1,\ldots,G^t=F'$$
such that, at odd-numbered steps, only odd-rank members may change, and at even-numbered steps, only even-rank members may change.
\end{lemma}

\begin{proposition}\label{prop:H}
For every prime power $q$ and every $k\ge1$, the graph $H_{k,q}$ is regular. Writing $\Delta_q$ for its degree, we have
\begin{equation}\label{eq:H-order}
  N_q:=|V(H_{k,q})|=\frac{[2k+1]_q!}{(q+1)^k},
\end{equation}
\begin{equation}\label{eq:H-degree}
  \Delta_q\le K_q:=(q+1)^k\left((q+1)^k-1\right),
\end{equation}
and
\begin{equation}\label{eq:H-diameter}
  \diam H_{k,q}=k.
\end{equation}
For fixed $k$, as $q\to\infty$ through prime powers,
\begin{equation}\label{eq:H-asymptotic}
  \Delta_q=(1+o(1))K_q,
  \qquad
  N_q=(1+o(1))\Delta_q^k.
\end{equation}
\end{proposition}

\begin{proof}
\smallskip
\noindent\emph{Order and degree.}
There are $k$ odd and $k$ even ranks among $1,\ldots,t-1$.  Fix
$P_{\mathrm{even}}\in\mathcal P_{\mathrm{even}}(V)$.  For each odd
$i\in\{1,\ldots,t-1\}$, a completion to a full flag must choose an
intermediate space
$$(P_{\mathrm{even}})_{i-1}<F_i<(P_{\mathrm{even}})_{i+1},$$
where $F_0=0$ and $F_t=V$ are understood.  The corresponding quotient has dimension two, and hence there are $q+1$ choices for $F_i$.  These choices are independent, so every even-rank partial flag has exactly $(q+1)^k$ completions.  The same argument applies to odd-rank partial flags.  Counting complete flags by their even-rank parts proves~\eqref{eq:H-order}.

Any two even-rank partial flags are carried to one another by an invertible linear map.  Since invertible linear maps preserve compatibility, they induce automorphisms of $H_{k,q}$.  Thus $H_{k,q}$ is vertex-transitive and hence regular.  From a fixed vertex there are exactly $(q+1)^k$ compatible odd-rank partial flags, and each of these is compatible with exactly $(q+1)^k-1$ other even-rank partial flags.  Different odd-rank partial flags may yield the same neighbor, so \eqref{eq:H-degree} follows.

\smallskip
\noindent\emph{Diameter.}
Take $P_{\mathrm{even}},P'_{\mathrm{even}}
\in\mathcal P_{\mathrm{even}}(V)$ and complete them to full flags $F$ and $F'$.  Apply Lemma~\ref{lem:routing} to obtain an odd-first route
$$F=G^0,G^1,\ldots,G^t=F'.$$
For $0\le j\le k$, set
$$P_j:=(G^{2j})_{\mathrm{even}}
      =(G^{2j+1})_{\mathrm{even}},$$
and, for $1\le j\le k$, set
$$Q_j:=(G^{2j-1})_{\mathrm{odd}}
      =(G^{2j})_{\mathrm{odd}}.$$
The equality for $P_j$ holds because odd steps leave all even-rank
members unchanged, while the equality for $Q_j$ holds because even
steps leave all odd-rank members unchanged. Thus $Q_j$ is compatible
with $P_{j-1}$ via $G^{2j-1}$ and with $P_j$ via $G^{2j}$. Hence
$$  P_{j-1} \text{ and } P_j \text{ are either equal or adjacent in } H_{k,q}.$$
Finally, since $G^0=F$ and $G^t=F'$, we have
$P_0=P_{\mathrm{even}}$ and $P_k=P'_{\mathrm{even}}$.
Deleting repetitions gives a path of length at most $k$.

For the reverse inequality, fix a basis $e_1,\ldots,e_t$.  For an odd- or even-rank partial flag $R=(R_i)$, put
$$\lambda(R):=\max\left(\{i:e_1\notin R_i\}\cup\{0\}\right).$$
Thus $\lambda(R)$ is the last recorded rank before $e_1$ appears.  If
$R\in\mathcal P_{\mathrm{even}}(V)$ and $S\in\mathcal P_{\mathrm{odd}}(V)$ are compatible, let $r$ be the first rank at which $e_1$ occurs in their common complete flag.  Then $\lambda(R)$ and $\lambda(S)$ are the largest even and odd ranks below $r$, respectively, so
$$|\lambda(R)-\lambda(S)|\le1.$$
It follows that adjacent vertices of $H_{k,q}$ have $\lambda$-values
differing by at most two.  For the standard and opposite complete flags
$$A_i=\operatorname{span}(e_1,\ldots,e_i),
  \qquad
  C_i=\operatorname{span}(e_t,e_{t-1},\ldots,e_{t+1-i}),$$
one has
$$\lambda(A_{\mathrm{even}})=0,
  \qquad
  \lambda(C_{\mathrm{even}})=2k.$$
A path of length $s$ between these vertices would therefore satisfy $2k\le2s$, so $s\ge k$.  This proves~\eqref{eq:H-diameter}.

\smallskip
\noindent\emph{Asymptotics.}
For fixed $k$, we have
$[2k+1]_q!
=q^{k(2k+1)}\left(1+O_k(q^{-1})\right).$
Equations~\eqref{eq:H-order} and~\eqref{eq:H-degree} give $N_q=q^{2k^2}\left(1+O_k(q^{-1})\right)$ and $K_q=q^{2k}\left(1+O_k(q^{-1})\right)$, 
and hence
\begin{equation}\label{eq:N-K-asymptotic}
  N_q=\left(1+O_k(q^{-1})\right)K_q^k.
\end{equation}
Since $N_q\to\infty$, the Moore bound forces $\Delta_q\to\infty$, and hence $N_q\le(1+o(1))\Delta_q^k$.
Together with~\eqref{eq:N-K-asymptotic} and $\Delta_q\le K_q$, this gives $\Delta_q\sim K_q$, and therefore $N_q\sim\Delta_q^k$.  This proves \eqref{eq:H-asymptotic}.
\end{proof}

\end{document}
